\subsection{A rank-one cache with an implicit query}
\label{sec:implicitrankone}

An explicit query can locate its score order.  When the query is supplied
only through fresh signs of an unknown scalar mean, that comparison is
unavailable.  A prefilled scalar factory still gives a polynomial source
bound at standard scaling.  We use a bias-corrected Bernstein expansion.
Bernstein approximation and smoothness-based simulation are established
Bernoulli-factory methods~\cite{nacuperes2005,holtz2011}; the role of
prefill here is to store their cache-dependent coefficient computations.
The correction is needed because the ordinary Bernstein approximation
has an error of order $m^{-1}$, whose direct mixture has infinite mean
source work.

\begin{lemma}[A summable Bernstein correction]
\label{lem:correctedbernstein}
Let $f\in C^4([0,1])$ satisfy $\|f\|_\infty\le4/5$.
Suppose $f$ and $f''$ have algorithms returning certified absolute
enclosures at every rational argument and requested precision.
Let known nonnegative rational numbers
$M_j\ge\|f^{(j)}\|_\infty$ be supplied for $j=2,3,4$.  Put
\[
 A=M_2+M_3+M_4,\qquad h(p)=p(1-p)f''(p),\qquad
 g_m(p)=f(p)-\frac{h(p)}{2m}.
\]
For $K\in\{0,\ldots,2m\}$, let $J$ be hypergeometric, counting
the successes in a uniformly selected half of a population of $2m$
bits containing $K$ successes.  Write $p=K/(2m)$ and
\[
 d_{m,K}=g_{2m}(p)-\mathbb E g_m(J/m).
\]
Then
\begin{equation}
 |d_{m,K}|\le\frac{A}{4m^2}.
 \label{eq:correctedcoefficientbound}
\end{equation}
If $m_0$ is a power of two satisfying
$m_0\ge\max\{1,2M_2,2\sqrt A\}$, there is an exact factory
returning a sign of mean $f(p)$ from independent source bits of
unknown mean $p$, using at most
\begin{equation}
 \frac78m_0+\frac{A}{m_0}
 \label{eq:correctedsourcecost}
\end{equation}
source bits in expectation.  This count does not include the work
needed to evaluate the known coefficient coins.
\end{lemma}

\paragraph{Proof idea.}
The second-derivative correction cancels the leading Bernstein approximation error.  Degree elevation is described by a symmetric hypergeometric sample, so a fourth-moment estimate leaves summable correction coefficients.  A mixture of these corrections yields the exact scalar mean.

\begin{proof}
Put $Y=J/m$ and $\Delta=Y-p$.  A sample and its complement have
the same distribution, so $\Delta$ is symmetric.  In particular
$\mathbb E\Delta=\mathbb E\Delta^3=0$, and the usual two-indicator
calculation gives
\[
 \nu:=\mathbb E\Delta^2=\frac{p(1-p)}{2m-1}.
\]
We also need the bound $\mathbb E\Delta^4\le m^{-2}$.
Here is an elementary derivation.  Reveal the $m$ sampled bits in
order, and take the Doob martingale of their final mean.  At reveal
$i$, its increment is
\[
 \frac{X_i-\mathbb E[X_i\mid X_1,\ldots,X_{i-1}]}{2m-i}.
\]
Its conditional range has length at most $1/m$.  The elementary
Hoeffding moment bound for a centered variable in an interval of
length $c$, $\mathbb E e^{tX}\le e^{t^2c^2/8}$, follows by
convexity of the exponential and one-variable differentiation.
Indeed convexity bounds the moment by that of the two endpoints
with the same mean.  After scaling by $c$, the logarithm of this
endpoint moment has value and derivative zero at the origin, and
its second derivative is the variance of a tilted Bernoulli bit,
at most $1/4$.  Integrating twice proves the bound for either sign
of $t$.
Iterating this bound gives
$\Pr(|\Delta|>x)\le2e^{-2mx^2}$.
Integrating the tail yields
$\mathbb E\Delta^4\le8\int_0^\infty x^3e^{-2mx^2}\,dx=m^{-2}$.

\paragraph{The coefficient bound.}
The second derivative of $h$ satisfies
\[
 H_2:=\|h''\|_\infty\le2M_2+2M_3+M_4/4.
\]
Taylor expansion of $f$ through degree three, symmetry, and the
fourth-moment bound give
\[
 \mathbb E f(Y)=f(p)+\frac12 f''(p)\nu+R,
 \qquad |R|\le\frac{M_4}{24m^2}.
\]
The first-order expansion of $h$ gives
$|\mathbb E h(Y)-h(p)|\le H_2\nu/2$.
The leading terms cancel as follows:
\[
 \frac{h(p)}{4m}-\frac12f''(p)\nu
 =-\frac{h(p)}{4m(2m-1)}.
\]
Since $|h|\le M_2/4$ and $\nu\le1/(4m)$, it follows that
\[
 |d_{m,K}|
 \le\frac1{m^2}\left(\frac{M_2+H_2}{16}+\frac{M_4}{24}\right)
 \le\frac{M_2+M_3+M_4}{4m^2}.
\]

\paragraph{The factory.}
Let $B_m g$ denote the degree-$m$ Bernstein polynomial of $g$.
Conditional on the total number $K$ of successes among $2m$
independent source bits, the number in their first half is the
hypergeometric variable above.  Therefore the coefficient of
degree $2m$ in $B_{2m}g_{2m}-B_mg_m$ is exactly $d_{m,K}$.
Also $B_mg_m\to f$ uniformly: ordinary Bernstein approximation
converges uniformly for continuous functions, and
$\|g_m-f\|_\infty\le M_2/(8m)$.
For the former claim, if $X\sim\operatorname{Bin}(m,p)$, split
$\mathbb E|f(X/m)-f(p)|$ at $|X/m-p|\le\delta$.
Uniform continuity bounds the first part, and Chebyshev's
inequality bounds the other part by
$2\|f\|_\infty/(4m\delta^2)$, uniformly in $p$.
Thus, for $m_\ell=2^\ell m_0$,
\[
 f=B_{m_0}g_{m_0}
   +\sum_{\ell\ge0}(B_{2m_\ell}g_{2m_\ell}-B_{m_\ell}g_{m_\ell}).
\]
Set $b_0=7/8$ and $c_\ell=A/(4m_\ell^2)$.
We have $\|g_{m_0}\|_\infty\le4/5+1/16<7/8$, and
\[
 \sum_{\ell\ge0}c_\ell=\frac{A}{3m_0^2}\le\frac1{12}.
\]
Select the base with probability $b_0$, select correction $\ell$
with probability $c_\ell$, and otherwise return a fair sign.
On the base draw $m_0$ source bits and, conditional on their count
$K$, return a known sign of mean $g_{m_0}(K/m_0)/b_0$.
On correction $\ell$ draw $2m_\ell$ bits and return a known sign
of mean $d_{m_\ell,K}/c_\ell$.
All these means are in $[-1,1]$.  The absolutely convergent
displayed identity proves exactness.  The mean source count is
\[
 b_0m_0+\sum_{\ell\ge0}2m_\ell c_\ell
 =\frac78m_0+\frac{A}{m_0}.
\]
If $A=0$, omit all corrections.  Otherwise the conditional level
law is $(3/4)4^{-\ell}$, sampled by pairs of fair random bits.
The level, source count, and every known coefficient comparison
are finite almost surely.
\end{proof}

\begin{lemma}[Uniform scalar softmax derivatives]
\label{lem:scalarsoftmaxderivatives}
Let $|v_j|\le1$ and $|a_j|\le\Lambda$, and define
\[
 u(p)=\frac{\sum_{j=1}^T v_j e^{a_jp+c_j}}
            {\sum_{j=1}^T e^{a_jp+c_j}},\qquad
 f(p)=u(p)/2.
\]
For arbitrary real offsets $c_j$ and $p\in[0,1]$,
\[
 \|f\|_\infty\le1/2,\qquad
 \|f''\|_\infty\le3\Lambda^2,\quad
 \|f'''\|_\infty\le13\Lambda^3,\quad
 \|f''''\|_\infty\le75\Lambda^4.
\]
The bounds are independent of $T$ and of the score offsets.
\end{lemma}

\paragraph{Proof idea.}
Differentiate the numerator identity repeatedly, dividing throughout by the positive softmax denominator.  The bounded slopes control each normalized derivative independently of the number of keys and their offsets.

\begin{proof}
Write the numerator as $N$ and denominator as $Z$.  Then
$|N^{(k)}|/Z\le\Lambda^k$ and
$|Z^{(k)}|/Z\le\Lambda^k$.
Differentiating $uZ=N$ inductively gives
\[
 |u'|\le2\Lambda,\quad |u''|\le6\Lambda^2,\quad
 |u'''|\le26\Lambda^3,\quad |u''''|\le150\Lambda^4.
\]
Dividing these inequalities by two proves the claim.  The radius
bound follows from $|u|\le1$.
\end{proof}

\begin{theorem}[A prefilled scalar softmax factory]
\label{thm:implicitrankone}
Consider the family in Lemma~\ref{lem:scalarsoftmaxderivatives}.
Assume $a_j,c_j,v_j$ and a rational bound $\Lambda$ have at most
$b$ bits, or that the slopes and offsets are rational multiples
of a common $1/\sqrt n$ with rational encodings of at most $b$ bits.
The online input consists of independent source bits of an unknown
mean $p\in[0,1]$.  Put
\[
 M_2=3\Lambda^2,\quad M_3=13\Lambda^3,\quad
 M_4=75\Lambda^4,\quad A=M_2+M_3+M_4,
\]
and let $m_0$ be the least power of two at least
$\max\{1,2M_2,2\sqrt A\}$.  Then a finite-bit data structure
samples a sign of mean $u(p)/2$ exactly, using at most
$(7/8)m_0+A/m_0=O(1+\Lambda^2)$ source requests and expected
auxiliary bit work
\[
 O\bigl((1+\Lambda^2)B^4\bigr),\qquad
 B=16+b+\lceil\log_2(T+n+\Lambda+2)\rceil.
\]
The preprocessing work is
$O(T^3(1+\Lambda^4)B^4)$ and stored bits are
$O(T(1+\Lambda^2)B)$.
All constants are absolute.  No exact-real operations are required.  The output is an attention-coordinate sign at radius
two.  It need not reveal an attention index.
More explicitly, the source count is less than
$45\max\{1,\Lambda^2\}$: one has
$A\le91\max\{\Lambda^2,\Lambda^4\}$,
$m_0<40\max\{1,\Lambda^2\}$, and $A/m_0\le m_0/4$.
\end{theorem}

\paragraph{Proof idea.}
Prefill stores the corrected Bernstein coefficients through a context-dependent cutoff.  The geometric correction law makes higher levels rare enough to evaluate when requested, and certified intervals resolve coefficient coins without altering their biases.

\begin{proof}
If $A=0$, omit every correction table and correction branch below;
the base branch alone suffices.  Assume $A>0$ whenever a correction
is normalized by its mixture mass.
We first specify a certified coefficient routine.  For rational
$x\in[0,1]$, a precision-$p$ enclosure for $f(x)$ and $f''(x)$
costs $O(T(B+\log m+p)^4)$ bit operations when
$x$ has $O(\log m)$ bits.  To see this, subtract the largest score
before evaluating exponentials.  Its weight is one, so the
denominator is at least one.  Weighted moments of the slopes up
to order two give $u,u',u''$ by the quotient rule, and division
by two gives $f,f''$.  The slope bounds require only
$O(B)$ extra guard bits.  Certified negative-exponential enclosures
have fourth-power bit cost by the preceding arithmetic lemma.
In the common quadratic-field case, score comparisons
reduce to rational comparisons, and certified square-root
enclosures suffice for all evaluations.

\paragraph{Correction coefficients.}
For a correction coefficient, the hypergeometric mass is proportional to
\[
 w_j=\binom{m}{j}\binom{m}{K-j},
 \quad \max(0,K-m)\le j\le\min(m,K).
\]
Its mode is $j_0=\lfloor K/2\rfloor$, possibly one of two tied
modes.  Set the mode weight to one and generate weights outwards.
The forward ratio is
\[
 \frac{w_{j+1}}{w_j}
 =\frac{(m-j)(K-j)}{(j+1)(m-K+j+1)};
\]
on the upper side it is at most one, and the reverse ratio is at
most one on the lower side.  Every ratio has $O(\log(m+2))$ bits.
Using fixed-point intervals at $r$ fractional bits, each outward
path accumulates at most $O(m2^{-r})$ error.  The total mass error
is $O(m^2 2^{-r})$.  The exact total is at least one because the
mode weight is one, so division is stable.  Taking
$r=p+2\lceil\log_2(m+2)\rceil+O(1)$ therefore gives the weighted
average to error $O(2^{-p})$.  Evaluate the $O(m)$ required
$g_m(j/m)$ values by the preceding routine with
$p+\lceil\log_2(m+2)\rceil+O(1)$ precision bits.
Thus a correction
coefficient has a precision-$p$ enclosure in
\[
 O(mT(B+\log(m+2)+p)^4)
\]
bit operations.  Normalizing by $c_\ell=A/(4m^2)$ adds only
$O(B+\log(m+2))$ precision bits.  No binomial integer
with $\Theta(m)$ bits is ever formed.

\paragraph{The prefilled table.}
Let $M$ be the smallest level $m_\ell$ at least $Tm_0$; then
$Tm_0\le M<2Tm_0$.  During prefill, store certified enclosures for
all base and correction coefficient coin probabilities through
level $M$.  For degree $2m$ use width at most $2^{-P_m}$, where
\[
 P_m=8+\lceil\log_2(Tm+2)\rceil.
\]
These dyadic interval data are used for exact comparisons with the
desired probabilities.  The sum of the table sizes is
$O(MB)$.  Evaluating all entries costs
$O(TM^2B^4)$, giving the stated bounds.

\paragraph{Online sampling.}
Online, sample the exact mixture of Lemma~\ref{lem:correctedbernstein}
and draw its required source bits.  Lemma~\ref{lem:newexactification}
decides each stored coefficient coin exactly.  Its fast enclosure is
the stored interval, of width $2^{-P_m}$, and its fallback is the
certified routine above, of cost $O(mT(B+\log(m+2)+t)^4)$ at
precision $t$.  Since $2^{P_m}\ge2^8(Tm+2)$ and $\log m=O(B)$ for
stored levels, the lemma bounds the expected work by $O(B^4)$.

\paragraph{Levels above the table.}
For a level $m>M$, compute the coefficient directly.  Its expected
cost is $O(mT(B+\log m)^4)$, and its mixture
probability is $A/(4m^2)$.  Summing over dyadic levels gives
\[
 \sum_{m>M}\frac{A}{4m^2}\,
    mT(B+\log m)^4
 =O\left(\frac{AT}{M}(B+\log M)^4\right)
 =O\left(\frac{A}{m_0}B^4\right).
\]
Source counting, source addresses, and mixture selection have the
same bound after logarithmic factors: the weighted level tail is
geometric.  Together with the base work and the displayed source
count, this proves the online bound.  Uniform comparisons stop
almost surely, and the infinite mixture assigns probability one
to finite levels.  The coefficient intervals only decide exact
coins, so the uniform Bernstein limit proves exactness for every
source bias, including zero and one.
\end{proof}

For a rank-one cache $k_j=c+t_j u$, choose a known $R$ with
$|\langle q,u\rangle|\le R$ on the allowed query domain, and
suppose fresh signs of mean $z=\langle q,u\rangle/R$ are available.
Writing $p=(1+z)/2$, the common cache offset cancels and the scores
have slopes $a_j=2\beta R t_j/n$ and offsets $c_j=-\beta R t_j/n$.
The theorem therefore applies with
$\Lambda\ge2\beta R\max_j|t_j|/n$.
For instance, if $k_j=c+t_j u$ with $|t_j|\le1$ and
$\|q\|_\infty,\|u\|_\infty\le1$, take $R=n$ and
any rational $\Lambda\ge2\beta$.  For $\beta\ge1$ one may choose
$\Lambda\le4\beta$.  One source sign of mean $z$ uses at most one
query-coordinate sign by a signed mixture.  At standard scaling
$\beta=\sqrt n$, the expected source and auxiliary costs are
$O(nB^4)$, independently of $T$ after the stated
polynomial prefill.  As after Theorem~\ref{thm:newrankone}, a unit
tanh row and a bounded two-logit head consume the radius-two output
through constant-cost scalar factories, so the exact final token has
the same asymptotic bound.  Tables can also be built for every cached
value coordinate; this multiplies the prefill and storage by the value
dimension.  The two-key lower bound of
Theorem~\ref{thm:rankonequerylaw} has order $\beta^2$ before the
finite-input cap.  It is a worst-case statement over caches.
